% --- Archon physics formalization source begin ---
% archon:physics
% archon:covers PhyXMiniProblems/problem_phyx_mini_0964.lean
% archon:source-report reports/phyx_mini/problem_phyx_mini_0964.source.json

\chapter{Physics problem phyx\_mini\_0964}
\label{ch:PhyXMiniProblems_problem_phyx_mini_0964}

\paragraph{Problem source.}
\#\# Physical scenario

Figure shows an end view of two long, parallel wires perpendicular to the \textbackslash{}( xy \textbackslash{})-plane, each carrying a current \textbackslash{}( I \textbackslash{}) but in opposite directions.

\#\# Question

What is the magnitude of \textbackslash{}( \textbackslash{}vec\{B\} \textbackslash{}) when \textbackslash{}( x \textbackslash{}gg a \textbackslash{})?

\#\# Answer choices

- A: A: \textbackslash{}frac\{\textbackslash{}mu\_0 I\}\{2\textbackslash{}pi x\}

- B: B: \textbackslash{}frac\{\textbackslash{}mu\_0 I a\}\{\textbackslash{}pi x\textasciicircum{}2\}

- C: C: \textbackslash{}frac\{\textbackslash{}mu\_0 I a\textasciicircum{}2\}\{\textbackslash{}pi x\textasciicircum{}3\}

- D: D: \textbackslash{}frac\{\textbackslash{}mu\_0 I\}\{\textbackslash{}pi(x\textasciicircum{}2 + a\textasciicircum{}2)\}

\#\# Figure caption (auxiliary; use the image as primary evidence)

The image depicts a two-dimensional coordinate system with axes labeled \textbackslash{}(x\textbackslash{}) and \textbackslash{}(y\textbackslash{}). There are two main points of interest marked on the vertical \textbackslash{}(y\textbackslash{})-axis:

1. A purple circle with a dot at its center labeled \textbackslash{}(I\textbackslash{}) resides at a distance \textbackslash{}(a\textbackslash{}) above the origin.
2. A purple circle with a cross in it, also labeled \textbackslash{}(I\textbackslash{}), is located a distance \textbackslash{}(a\textbackslash{}) below the origin.

A black point labeled \textbackslash{}(P\textbackslash{}) is marked on the positive \textbackslash{}(x\textbackslash{})-axis. The horizontal distance from the origin \textbackslash{}(O\textbackslash{}) to point \textbackslash{}(P\textbackslash{}) is denoted by \textbackslash{}(x\textbackslash{}).

The image includes arrows indicating these distances:

- Vertical arrows indicate the distances \textbackslash{}(a\textbackslash{}) from the origin to each \textbackslash{}(I\textbackslash{}) point on the \textbackslash{}(y\textbackslash{})-axis.
- A horizontal arrow indicates the distance \textbackslash{}(x\textbackslash{}) from the origin \textbackslash{}(O\textbackslash{}) to point \textbackslash{}(P\textbackslash{}) on the \textbackslash{}(x\textbackslash{})-axis.

\#\# Dataset metadata

Magnetism; Physical Model Grounding Reasoning, Multi-Formula Reasoning

\paragraph{Recorded answer/context.}
B: B: \textbackslash{}frac\{\textbackslash{}mu\_0 I a\}\{\textbackslash{}pi x\textasciicircum{}2\}

\paragraph{Figure/image path.}
/root/proposal\_for\_physic/science-mango/phyx\_mini\_run/phyx\_data/test\_image/964.png

\paragraph{Formalization target.}
create a compiling Lean file with sorry bodies at `PhyXMiniProblems/problem\_phyx\_mini\_0964.lean`. The Lean declarations must preserve the physical quantities, dimensions or dimensional roles, figure labels, governing-law hypotheses, and final relation expressed by this problem.
Use Mathlib/Physlib names found through LeanExplore where available. If a physics API is missing, introduce faithful local abstractions rather than scalar placeholder aliases.

\begin{theorem}[Physics formalization target]
\label{thm:physics:phyx_mini_0964:target}
\lean{PhyXMiniProblems.ProblemPhyXMini0964.magneticFieldMagnitude_farField_eq_answer_B}
\uses{def:physics:phyx-mini-0964:phyxminiproblems-problemphyxmini0964-lengthinmeters, def:physics:phyx-mini-0964:phyxminiproblems-problemphyxmini0964-currentinamperes, def:physics:phyx-mini-0964:phyxminiproblems-problemphyxmini0964-permeabilityinteslametersperampere, def:physics:phyx-mini-0964:phyxminiproblems-problemphyxmini0964-opposedparallelwiresetup, def:physics:phyx-mini-0964:phyxminiproblems-problemphyxmini0964-magneticfieldmagnitudeonxaxisinteslas, def:physics:phyx-mini-0964:phyxminiproblems-problemphyxmini0964-matchesopposedparallelwirescenario, def:physics:phyx-mini-0964:phyxminiproblems-problemphyxmini0964-matchesprimaryopposedwirefigure, def:physics:phyx-mini-0964:phyxminiproblems-problemphyxmini0964-satisfiesopposedwiremagneticfieldlaws}
The assigned autoformalize agent should translate this physics problem into Lean declarations in the covered file, with theorem and lemma proof bodies written as `by sorry`.
\end{theorem}
\begin{proof}
This is an autoformalization task, not a proof task. Produce statements that can later be proved by the physics prover without weakening the physical model.
\end{proof}

% --- Archon declaration topology begin ---
\section{Declaration topology}
\begin{definition}[electric Current Dimension]
  \label{def:physics:phyx-mini-0964:phyxminiproblems-problemphyxmini0964-electriccurrentdimension}
  \lean{PhyXMiniProblems.ProblemPhyXMini0964.electricCurrentDimension}
  Electric current has physical dimension C T⁻¹
\end{definition}
\begin{proof}
  This is the declared model definition of electric Current Dimension.
\end{proof}
\begin{definition}[magnetic Permeability Dimension]
  \label{def:physics:phyx-mini-0964:phyxminiproblems-problemphyxmini0964-magneticpermeabilitydimension}
  \lean{PhyXMiniProblems.ProblemPhyXMini0964.magneticPermeabilityDimension}
  Magnetic permeability has SI dimension M L C⁻²
\end{definition}
\begin{proof}
  This is the declared model definition of magnetic Permeability Dimension.
\end{proof}
\begin{definition}[magnetic Flux Density Dimension]
  \label{def:physics:phyx-mini-0964:phyxminiproblems-problemphyxmini0964-magneticfluxdensitydimension}
  \lean{PhyXMiniProblems.ProblemPhyXMini0964.magneticFluxDensityDimension}
  Magnetic flux density has SI dimension M T⁻¹ C⁻¹
\end{definition}
\begin{proof}
  This is the declared model definition of magnetic Flux Density Dimension.
\end{proof}
\begin{definition}[Length Quantity]
  \label{def:physics:phyx-mini-0964:phyxminiproblems-problemphyxmini0964-lengthquantity}
  \lean{PhyXMiniProblems.ProblemPhyXMini0964.LengthQuantity}
  A nonnegative, unit-independent physical length
\end{definition}
\begin{proof}
  This is the declared model definition of Length Quantity.
\end{proof}
\begin{definition}[Electric Current Magnitude]
  \label{def:physics:phyx-mini-0964:phyxminiproblems-problemphyxmini0964-electriccurrentmagnitude}
  \lean{PhyXMiniProblems.ProblemPhyXMini0964.ElectricCurrentMagnitude}
  \uses{def:physics:phyx-mini-0964:phyxminiproblems-problemphyxmini0964-electriccurrentdimension}
  A nonnegative, unit-independent electric-current magnitude
\end{definition}
\begin{proof}
  This is the declared model definition of Electric Current Magnitude.
\end{proof}
\begin{definition}[Magnetic Permeability Quantity]
  \label{def:physics:phyx-mini-0964:phyxminiproblems-problemphyxmini0964-magneticpermeabilityquantity}
  \lean{PhyXMiniProblems.ProblemPhyXMini0964.MagneticPermeabilityQuantity}
  \uses{def:physics:phyx-mini-0964:phyxminiproblems-problemphyxmini0964-magneticpermeabilitydimension}
  A nonnegative, unit-independent magnetic permeability
\end{definition}
\begin{proof}
  This is the declared model definition of Magnetic Permeability Quantity.
\end{proof}
\begin{definition}[nonnegative SIReadout]
  \label{def:physics:phyx-mini-0964:phyxminiproblems-problemphyxmini0964-nonnegativesireadout}
  \lean{PhyXMiniProblems.ProblemPhyXMini0964.nonnegativeSIReadout}
  Read a nonnegative dimensionful quantity in coherent SI units
\end{definition}
\begin{proof}
  This is the declared model definition of nonnegative SIReadout.
\end{proof}
\begin{definition}[length In Meters]
  \label{def:physics:phyx-mini-0964:phyxminiproblems-problemphyxmini0964-lengthinmeters}
  \lean{PhyXMiniProblems.ProblemPhyXMini0964.lengthInMeters}
  \uses{def:physics:phyx-mini-0964:phyxminiproblems-problemphyxmini0964-lengthquantity, def:physics:phyx-mini-0964:phyxminiproblems-problemphyxmini0964-nonnegativesireadout}
  Read a physical length in metres
\end{definition}
\begin{proof}
  This is the declared model definition of length In Meters.
\end{proof}
\begin{definition}[current In Amperes]
  \label{def:physics:phyx-mini-0964:phyxminiproblems-problemphyxmini0964-currentinamperes}
  \lean{PhyXMiniProblems.ProblemPhyXMini0964.currentInAmperes}
  \uses{def:physics:phyx-mini-0964:phyxminiproblems-problemphyxmini0964-electriccurrentmagnitude, def:physics:phyx-mini-0964:phyxminiproblems-problemphyxmini0964-nonnegativesireadout}
  Read an electric-current magnitude in amperes
\end{definition}
\begin{proof}
  This is the declared model definition of current In Amperes.
\end{proof}
\begin{definition}[permeability In Tesla Meters Per Ampere]
  \label{def:physics:phyx-mini-0964:phyxminiproblems-problemphyxmini0964-permeabilityinteslametersperampere}
  \lean{PhyXMiniProblems.ProblemPhyXMini0964.permeabilityInTeslaMetersPerAmpere}
  \uses{def:physics:phyx-mini-0964:phyxminiproblems-problemphyxmini0964-magneticpermeabilityquantity, def:physics:phyx-mini-0964:phyxminiproblems-problemphyxmini0964-nonnegativesireadout}
  Read permeability in tesla-metres per ampere
\end{definition}
\begin{proof}
  This is the declared model definition of permeability In Tesla Meters Per Ampere.
\end{proof}
\begin{definition}[Spatial Vector]
  \label{def:physics:phyx-mini-0964:phyxminiproblems-problemphyxmini0964-spatialvector}
  \lean{PhyXMiniProblems.ProblemPhyXMini0964.SpatialVector}
  A three-dimensional spatial vector used for directions and SI components
\end{definition}
\begin{proof}
  This is the declared model definition of Spatial Vector.
\end{proof}
\begin{definition}[Coordinate Axis]
  \label{def:physics:phyx-mini-0964:phyxminiproblems-problemphyxmini0964-coordinateaxis}
  \lean{PhyXMiniProblems.ProblemPhyXMini0964.CoordinateAxis}
  The coordinate axes; only x and y appear in the end-view raster
\end{definition}
\begin{proof}
  This is the declared model definition of Coordinate Axis.
\end{proof}
\begin{definition}[axis Direction]
  \label{def:physics:phyx-mini-0964:phyxminiproblems-problemphyxmini0964-axisdirection}
  \lean{PhyXMiniProblems.ProblemPhyXMini0964.axisDirection}
  \uses{def:physics:phyx-mini-0964:phyxminiproblems-problemphyxmini0964-spatialvector, def:physics:phyx-mini-0964:phyxminiproblems-problemphyxmini0964-coordinateaxis}
  Unit vector associated with a coordinate axis
\end{definition}
\begin{proof}
  This is the declared model definition of axis Direction.
\end{proof}
\begin{definition}[spatial Cross Product]
  \label{def:physics:phyx-mini-0964:phyxminiproblems-problemphyxmini0964-spatialcrossproduct}
  \lean{PhyXMiniProblems.ProblemPhyXMini0964.spatialCrossProduct}
  \uses{def:physics:phyx-mini-0964:phyxminiproblems-problemphyxmini0964-spatialvector}
  Mathlib's coordinate cross product transported to the Euclidean-space type used by Physlib magnetic fields
\end{definition}
\begin{proof}
  This is the declared model definition of spatial Cross Product.
\end{proof}
\begin{definition}[point On XAxis Meters]
  \label{def:physics:phyx-mini-0964:phyxminiproblems-problemphyxmini0964-pointonxaxismeters}
  \lean{PhyXMiniProblems.ProblemPhyXMini0964.pointOnXAxisMeters}
  \uses{def:physics:phyx-mini-0964:phyxminiproblems-problemphyxmini0964-axisdirection}
  The spatial point with coherent-SI coordinates (x, 0, 0)
\end{definition}
\begin{proof}
  This is the declared model definition of point On XAxis Meters.
\end{proof}
\begin{definition}[Wire]
  \label{def:physics:phyx-mini-0964:phyxminiproblems-problemphyxmini0964-wire}
  \lean{PhyXMiniProblems.ProblemPhyXMini0964.Wire}
  The upper and lower long wires in the primary figure
\end{definition}
\begin{proof}
  This is the declared model definition of Wire.
\end{proof}
\begin{definition}[Wire Model]
  \label{def:physics:phyx-mini-0964:phyxminiproblems-problemphyxmini0964-wiremodel}
  \lean{PhyXMiniProblems.ProblemPhyXMini0964.WireModel}
  Idealization assigned to each current-carrying source
\end{definition}
\begin{proof}
  This is the declared model definition of Wire Model.
\end{proof}
\begin{definition}[Current Glyph]
  \label{def:physics:phyx-mini-0964:phyxminiproblems-problemphyxmini0964-currentglyph}
  \lean{PhyXMiniProblems.ProblemPhyXMini0964.CurrentGlyph}
  End-view glyph for a page-normal current direction
\end{definition}
\begin{proof}
  This is the declared model definition of Current Glyph.
\end{proof}
\begin{definition}[Figure Point]
  \label{def:physics:phyx-mini-0964:phyxminiproblems-problemphyxmini0964-figurepoint}
  \lean{PhyXMiniProblems.ProblemPhyXMini0964.FigurePoint}
  Labeled points and wire centers visible in image 964.png
\end{definition}
\begin{proof}
  This is the declared model definition of Figure Point.
\end{proof}
\begin{definition}[Distance Guide]
  \label{def:physics:phyx-mini-0964:phyxminiproblems-problemphyxmini0964-distanceguide}
  \lean{PhyXMiniProblems.ProblemPhyXMini0964.DistanceGuide}
  The three double-arrow distance guides shown in the raster
\end{definition}
\begin{proof}
  This is the declared model definition of Distance Guide.
\end{proof}
\begin{definition}[expected Distance Symbol]
  \label{def:physics:phyx-mini-0964:phyxminiproblems-problemphyxmini0964-expecteddistancesymbol}
  \lean{PhyXMiniProblems.ProblemPhyXMini0964.expectedDistanceSymbol}
  \uses{def:physics:phyx-mini-0964:phyxminiproblems-problemphyxmini0964-distanceguide}
  Printed symbol beside each distance guide
\end{definition}
\begin{proof}
  This is the declared model definition of expected Distance Symbol.
\end{proof}
\begin{definition}[Opposed Parallel Wire Figure]
  \label{def:physics:phyx-mini-0964:phyxminiproblems-problemphyxmini0964-opposedparallelwirefigure}
  \lean{PhyXMiniProblems.ProblemPhyXMini0964.OpposedParallelWireFigure}
  \uses{def:physics:phyx-mini-0964:phyxminiproblems-problemphyxmini0964-coordinateaxis, def:physics:phyx-mini-0964:phyxminiproblems-problemphyxmini0964-wire, def:physics:phyx-mini-0964:phyxminiproblems-problemphyxmini0964-currentglyph, def:physics:phyx-mini-0964:phyxminiproblems-problemphyxmini0964-figurepoint, def:physics:phyx-mini-0964:phyxminiproblems-problemphyxmini0964-distanceguide}
  Literal presentation data from the primary raster. Coordinate values are coherent-SI metre readouts; the physical lengths that calibrate them live in OpposedParallelWireSetup
\end{definition}
\begin{proof}
  This is the declared model definition of Opposed Parallel Wire Figure.
\end{proof}
\begin{definition}[Opposed Parallel Wire Setup]
  \label{def:physics:phyx-mini-0964:phyxminiproblems-problemphyxmini0964-opposedparallelwiresetup}
  \lean{PhyXMiniProblems.ProblemPhyXMini0964.OpposedParallelWireSetup}
  \uses{def:physics:phyx-mini-0964:phyxminiproblems-problemphyxmini0964-lengthquantity, def:physics:phyx-mini-0964:phyxminiproblems-problemphyxmini0964-electriccurrentmagnitude, def:physics:phyx-mini-0964:phyxminiproblems-problemphyxmini0964-magneticpermeabilityquantity, def:physics:phyx-mini-0964:phyxminiproblems-problemphyxmini0964-spatialvector, def:physics:phyx-mini-0964:phyxminiproblems-problemphyxmini0964-coordinateaxis, def:physics:phyx-mini-0964:phyxminiproblems-problemphyxmini0964-wire, def:physics:phyx-mini-0964:phyxminiproblems-problemphyxmini0964-wiremodel, def:physics:phyx-mini-0964:phyxminiproblems-problemphyxmini0964-opposedparallelwirefigure}
  Independent physical quantities, source fields, and the total observable field. None of these fields is defined from an answer formula
\end{definition}
\begin{proof}
  This is the declared model definition of Opposed Parallel Wire Setup.
\end{proof}
\begin{definition}[magnetic Field Vector Due To Wire On XAxis In Teslas]
  \label{def:physics:phyx-mini-0964:phyxminiproblems-problemphyxmini0964-magneticfieldvectorduetowireonxaxisinteslas}
  \lean{PhyXMiniProblems.ProblemPhyXMini0964.magneticFieldVectorDueToWireOnXAxisInTeslas}
  \uses{def:physics:phyx-mini-0964:phyxminiproblems-problemphyxmini0964-spatialvector, def:physics:phyx-mini-0964:phyxminiproblems-problemphyxmini0964-pointonxaxismeters, def:physics:phyx-mini-0964:phyxminiproblems-problemphyxmini0964-wire, def:physics:phyx-mini-0964:phyxminiproblems-problemphyxmini0964-opposedparallelwiresetup}
  The field due to one wire at (x,0,0), read as a vector in teslas
\end{definition}
\begin{proof}
  This is the declared model definition of magnetic Field Vector Due To Wire On XAxis In Teslas.
\end{proof}
\begin{definition}[magnetic Field Vector On XAxis In Teslas]
  \label{def:physics:phyx-mini-0964:phyxminiproblems-problemphyxmini0964-magneticfieldvectoronxaxisinteslas}
  \lean{PhyXMiniProblems.ProblemPhyXMini0964.magneticFieldVectorOnXAxisInTeslas}
  \uses{def:physics:phyx-mini-0964:phyxminiproblems-problemphyxmini0964-spatialvector, def:physics:phyx-mini-0964:phyxminiproblems-problemphyxmini0964-pointonxaxismeters, def:physics:phyx-mini-0964:phyxminiproblems-problemphyxmini0964-opposedparallelwiresetup}
  The total magnetic-field vector at (x,0,0), in teslas
\end{definition}
\begin{proof}
  This is the declared model definition of magnetic Field Vector On XAxis In Teslas.
\end{proof}
\begin{definition}[magnetic Field Magnitude On XAxis In Teslas]
  \label{def:physics:phyx-mini-0964:phyxminiproblems-problemphyxmini0964-magneticfieldmagnitudeonxaxisinteslas}
  \lean{PhyXMiniProblems.ProblemPhyXMini0964.magneticFieldMagnitudeOnXAxisInTeslas}
  \uses{def:physics:phyx-mini-0964:phyxminiproblems-problemphyxmini0964-opposedparallelwiresetup, def:physics:phyx-mini-0964:phyxminiproblems-problemphyxmini0964-magneticfieldvectoronxaxisinteslas}
  The magnitude of the total field at (x,0,0), in teslas
\end{definition}
\begin{proof}
  This is the declared model definition of magnetic Field Magnitude On XAxis In Teslas.
\end{proof}
\begin{definition}[Matches Opposed Parallel Wire Scenario]
  \label{def:physics:phyx-mini-0964:phyxminiproblems-problemphyxmini0964-matchesopposedparallelwirescenario}
  \lean{PhyXMiniProblems.ProblemPhyXMini0964.MatchesOpposedParallelWireScenario}
  \uses{def:physics:phyx-mini-0964:phyxminiproblems-problemphyxmini0964-lengthinmeters, def:physics:phyx-mini-0964:phyxminiproblems-problemphyxmini0964-currentinamperes, def:physics:phyx-mini-0964:phyxminiproblems-problemphyxmini0964-permeabilityinteslametersperampere, def:physics:phyx-mini-0964:phyxminiproblems-problemphyxmini0964-axisdirection, def:physics:phyx-mini-0964:phyxminiproblems-problemphyxmini0964-opposedparallelwiresetup}
  The prose-level idealizations, page-normal orientations, and positivity of the physical magnitudes. A single current-magnitude field records that both wires carry the same I; the signed unit directions record that the currents are opposite
\end{definition}
\begin{proof}
  This is the declared model definition of Matches Opposed Parallel Wire Scenario.
\end{proof}
\begin{definition}[Matches Primary Opposed Wire Figure]
  \label{def:physics:phyx-mini-0964:phyxminiproblems-problemphyxmini0964-matchesprimaryopposedwirefigure}
  \lean{PhyXMiniProblems.ProblemPhyXMini0964.MatchesPrimaryOpposedWireFigure}
  \uses{def:physics:phyx-mini-0964:phyxminiproblems-problemphyxmini0964-lengthinmeters, def:physics:phyx-mini-0964:phyxminiproblems-problemphyxmini0964-spatialvector, def:physics:phyx-mini-0964:phyxminiproblems-problemphyxmini0964-axisdirection, def:physics:phyx-mini-0964:phyxminiproblems-problemphyxmini0964-pointonxaxismeters, def:physics:phyx-mini-0964:phyxminiproblems-problemphyxmini0964-expecteddistancesymbol, def:physics:phyx-mini-0964:phyxminiproblems-problemphyxmini0964-opposedparallelwiresetup}
  The axes, symbols, dot/cross convention, guide endpoints, and metric coordinates transcribed from the supplied image. This fixes the wire centers at (0, ±a) and the displayed point P at (x,0) without asserting any magnetic-field value
\end{definition}
\begin{proof}
  This is the declared model definition of Matches Primary Opposed Wire Figure.
\end{proof}
\begin{definition}[Satisfies Opposed Wire Magnetic Field Laws]
  \label{def:physics:phyx-mini-0964:phyxminiproblems-problemphyxmini0964-satisfiesopposedwiremagneticfieldlaws}
  \lean{PhyXMiniProblems.ProblemPhyXMini0964.SatisfiesOpposedWireMagneticFieldLaws}
  \uses{def:physics:phyx-mini-0964:phyxminiproblems-problemphyxmini0964-currentinamperes, def:physics:phyx-mini-0964:phyxminiproblems-problemphyxmini0964-permeabilityinteslametersperampere, def:physics:phyx-mini-0964:phyxminiproblems-problemphyxmini0964-spatialcrossproduct, def:physics:phyx-mini-0964:phyxminiproblems-problemphyxmini0964-wire, def:physics:phyx-mini-0964:phyxminiproblems-problemphyxmini0964-opposedparallelwiresetup}
  The governing electromagnetic laws used in the intended derivation. For a wire with unit current direction k and displacement vector r from the wire to the observation point, the vector field is B = μ₀ I / (2 π ‖r‖²) • (k × r). This is the general infinite-straight-wire law, not the requested net or far-field result. The second field records ordinary magnetic superposition
\end{definition}
\begin{proof}
  This is the declared model definition of Satisfies Opposed Wire Magnetic Field Laws.
\end{proof}
\begin{lemma}[magnetic Field Vector On XAxis exact]
  \label{lem:physics:phyx-mini-0964:phyxminiproblems-problemphyxmini0964-magneticfieldvectoronxaxis-exact}
  \lean{PhyXMiniProblems.ProblemPhyXMini0964.magneticFieldVectorOnXAxis_exact}
  \uses{def:physics:phyx-mini-0964:phyxminiproblems-problemphyxmini0964-lengthinmeters, def:physics:phyx-mini-0964:phyxminiproblems-problemphyxmini0964-currentinamperes, def:physics:phyx-mini-0964:phyxminiproblems-problemphyxmini0964-permeabilityinteslametersperampere, def:physics:phyx-mini-0964:phyxminiproblems-problemphyxmini0964-axisdirection, def:physics:phyx-mini-0964:phyxminiproblems-problemphyxmini0964-opposedparallelwiresetup, def:physics:phyx-mini-0964:phyxminiproblems-problemphyxmini0964-magneticfieldvectoronxaxisinteslas, def:physics:phyx-mini-0964:phyxminiproblems-problemphyxmini0964-matchesopposedparallelwirescenario, def:physics:phyx-mini-0964:phyxminiproblems-problemphyxmini0964-matchesprimaryopposedwirefigure, def:physics:phyx-mini-0964:phyxminiproblems-problemphyxmini0964-satisfiesopposedwiremagneticfieldlaws}
  The right-hand rule and the symmetric geometry make the y components from the two wires cancel and the x components add. This exact vector identity is a derived result, not a governing-law premise
\end{lemma}
\begin{proof}
  The conclusion follows from the governing relations and figure readouts named in the statement of magnetic Field Vector On XAxis exact.
\end{proof}
\begin{lemma}[magnetic Field Magnitude On XAxis exact]
  \label{lem:physics:phyx-mini-0964:phyxminiproblems-problemphyxmini0964-magneticfieldmagnitudeonxaxis-exact}
  \lean{PhyXMiniProblems.ProblemPhyXMini0964.magneticFieldMagnitudeOnXAxis_exact}
  \uses{def:physics:phyx-mini-0964:phyxminiproblems-problemphyxmini0964-lengthinmeters, def:physics:phyx-mini-0964:phyxminiproblems-problemphyxmini0964-currentinamperes, def:physics:phyx-mini-0964:phyxminiproblems-problemphyxmini0964-permeabilityinteslametersperampere, def:physics:phyx-mini-0964:phyxminiproblems-problemphyxmini0964-opposedparallelwiresetup, def:physics:phyx-mini-0964:phyxminiproblems-problemphyxmini0964-magneticfieldmagnitudeonxaxisinteslas, def:physics:phyx-mini-0964:phyxminiproblems-problemphyxmini0964-matchesopposedparallelwirescenario, def:physics:phyx-mini-0964:phyxminiproblems-problemphyxmini0964-matchesprimaryopposedwirefigure, def:physics:phyx-mini-0964:phyxminiproblems-problemphyxmini0964-satisfiesopposedwiremagneticfieldlaws}
  Exact finite-distance magnitude before making the x ≫ a approximation
\end{lemma}
\begin{proof}
  The conclusion follows from the governing relations and figure readouts named in the statement of magnetic Field Magnitude On XAxis exact.
\end{proof}
\begin{definition}[Answer Choice]
  \label{def:physics:phyx-mini-0964:phyxminiproblems-problemphyxmini0964-answerchoice}
  \lean{PhyXMiniProblems.ProblemPhyXMini0964.AnswerChoice}
  Labels of the four formulas displayed in the multiple-choice question
\end{definition}
\begin{proof}
  This is the declared model definition of Answer Choice.
\end{proof}
\begin{definition}[displayed Far Field Formula In Teslas]
  \label{def:physics:phyx-mini-0964:phyxminiproblems-problemphyxmini0964-displayedfarfieldformulainteslas}
  \lean{PhyXMiniProblems.ProblemPhyXMini0964.displayedFarFieldFormulaInTeslas}
  \uses{def:physics:phyx-mini-0964:phyxminiproblems-problemphyxmini0964-lengthinmeters, def:physics:phyx-mini-0964:phyxminiproblems-problemphyxmini0964-currentinamperes, def:physics:phyx-mini-0964:phyxminiproblems-problemphyxmini0964-permeabilityinteslametersperampere, def:physics:phyx-mini-0964:phyxminiproblems-problemphyxmini0964-opposedparallelwiresetup, def:physics:phyx-mini-0964:phyxminiproblems-problemphyxmini0964-answerchoice}
  Scalar coherent-SI readings of the four printed formulas. Some distractors are dimensionally invalid as physical formulas; recording them as displayed scalar expressions does not promote them to governing laws
\end{definition}
\begin{proof}
  This is the declared model definition of displayed Far Field Formula In Teslas.
\end{proof}
% --- Archon declaration topology end ---
% --- Archon physics formalization source end ---
